\documentclass[11pt]{article}
\usepackage{amsmath,amssymb,hyperref}
\title{Exact Continuous Recovery in a Finite Union of Wide Corridor Cells}
\date{2026-11-27}
\begin{document}
\maketitle

\paragraph{Audit.}
The controlled variables are the ordered robot positions on a line,
with coordinate speed bounds \(v_i\geq0\). A closed, rational polyhedral
set of allowed positions is the same union of coverage cells as in
the preceding wide-corridor model. All admitted survivor sets have the
required sensing coverage of the protected interval, while the
clearance, coordinate and required radio span constraints hold throughout.
The admissible environment, loss family, track ordering, endpoints and
bounds stay fixed over the leg. The question is the minimum time
from \(A\) to \(B\), not whether an initially broken requirement can
be counted as satisfied on time zero. There is no imposed bound on
acceleration or jerk and no mandatory rest interval, no changes to the
admitted loss sets, and no private observations in this solve.

\paragraph{A finite exact reduction.}
Let \(P_1,\ldots,P_m\) be a finite closed convex cover of the allowable
space \(S\). Suppose a continuous speed-feasible trajectory \(x\)
runs through \(S\) for time \(T\). If \(T=0\), its endpoints coincide.
Otherwise choose any cell containing \(x(0)\) and let \(t_1\) be
the last time through \(T\) at which the path is in that cell.
The inverse image is closed, so the maximum is attained.
If \(t_1<T\), take a sequence of times strictly above \(t_1\)
converging to it. Finiteness gives a cell visited by infinitely
many of those times. Closedness puts \(x(t_1)\) in that cell too;
its last visit is strictly later. Switch there, and iterate.
Every selected cell has last visit strictly greater than those
already used, with the possible initial zero-length stay.
Thus the procedure terminates at \(T\) using at most \(m\) distinct
cells. It does not require a minimum dwell or a finite number of
crossings of the original trajectory.

Replace each portion between consecutive chosen cuts by its affine
chord. Both endpoints of this chord lie in that cell, so the
continuous chord stays inside it. For each robot the displacement
is at most \(v_i\) times the original duration. This proves
speed feasibility with exactly those durations, including fixed
coordinates and any zero-duration identical endpoints. At a seam
velocity may jump. This is why the reduction applies to positions
with box speed bounds, not to an acceleration-limited state
augmented by velocities. For a cell that already contains
\(B\), any feasible remaining path can similarly be replaced
by an in-cell chord to \(B\) using no more time.

For each possible simple cell route choose seam positions \(z_j\)
and durations \(d_j\geq0\). Both endpoints of segment \(j\) must
be in its cell and
\[
  -v_i d_j\le z_{j+1,i}-z_{j,i}\le v_i d_j .
\]
Solve the rational linear program minimizing \(\sum_j d_j\),
with the fixed start and finish positions. Nonintersecting cells
can be pruned before the enumeration. At fixed speed the feasible
chords concatenate, giving the converse direction. A rational LP
with a finite infimum here attains it; nonnegative duration and
bounded position variables allow the search for a minimizer
to be restricted to the objective sublevel of an arbitrary
feasible witness. Exact finite LP algorithms can therefore decide
each route, and the finite minimum is exact. No feasible route
implies no continuous speed-feasible path. The simple route
count can grow exponentially or worse; no online tractability or
performance gain over a stronger exact geometric solver is claimed.

\paragraph{Certificates, not a floating decision.}
The held-in harness uses the source-indexed cells from the earlier
replay. Intersection pruning uses exact shortest-path difference
constraints with a zero reference coordinate, not the floating
adjacency test. For each unpruned route HiGHS proposes a primal
and dual of the minimization form \(A x\le b\) with unrestricted
variables (positive durations represented as rows). The harness
reconstructs rational candidate values and checks
\[
  A x\le b,\qquad A^\top y=c,\quad y\le0,\qquad c^\top x=b^\top y .
\]
For an infeasibility claim it checks a rational \(y\geq0\) with
\(A^\top y=0\) and \(b^\top y=-1\). These are exact comparisons,
including objective equality, not nominal tolerances. The
harness stops unconfirmed on a failed reconstruction; it does not
implement general exact LP fallback. The conceptual completeness
claim uses an exact solver, not indefinite trust in rational
rounding of a floating estimate.

\paragraph{Checks against a matched direct planner.}
All 77 old ascending pairs, 24 additional wide pairs with
noncomparable endpoints, and four controls agree with the
previous uninstrumented same-information route planner. There
were 744 audited route LPs, three of them infeasible with
exact Farkas witnesses. The 119 winning chords passed the
independent exact survivor-interval-order test from the wide
corridor source. The two broad examples still take exactly 7 and
\(7/2\); the disconnected case previously misclassified by a
coarse sample remains infeasible. In a separate control the static
cell graph is connected, but a robot fixed by its zero speed cap
would have to move to attain the finish, so the route LP refuses
correctly. Refusal is not an otherwise successful service claim.

This is an exact speed-only specialization of the
\href{https://arxiv.org/abs/2205.04422}{graphs-of-convex-sets
motion-planning framework}, using the previous corridor cell
construction. It closes the continuous finite-route argument and
certificate-verification gap for that scoped wide model. It does
not establish that the sensors or radio physically implement
the specified inequalities, does not solve the acceleration-limited
or time-varying support problem, and is not an independent contract
acceptance test.
\end{document}
